\section{Evaluation}
\label{sec:evaluation}

The evaluation is organized around three questions tied to the paper's claims.

\paragraph{RQ1: What assurance is present in the artifact?}
The coverage report records, for every operation, the available numerical
theorem, condition bridge, C-body refinement, check proof, guarded-wrapper
proof, and linked interface.  It separately reports handwritten C, numerical
proof, VST proof, and generated Clight sizes; clean-build time; dependency
versions; and a transitive assumption audit for the linked library.  These
measurements establish coverage and reproducibility, not the scientific value
of the proof by line count alone.

\paragraph{RQ2: What does the checked interface cost?}
We compare guarded and unchecked functions with QD, matching algorithmic
variants where the reference library offers multiple choices~\cite{QD}.  The
unchecked core separates numerical work from guard cost.  The evaluation also
separates a CompCert-produced library from builds in which another compiler
inlines the same source, because those configurations have different
performance and verification boundaries.  Each table records the processor,
operating system, compiler and flags, QD version, input distribution, number of
trials, and summary statistic.

\paragraph{RQ3: What domains do the guards expose?}
Boundary-focused tests sweep representative values around overflow and
normal/subnormal transitions.  For each case, the report gives the input bit
patterns, check outcome, contract tier, and applicable theorem.  An MPFR oracle
is used to inspect experimental behavior and validate the harness; it is not
evidence for the machine-checked theorem.  Comparisons with QD are presented as
contract and domain comparisons rather than claims that QD violates a contract
it does not advertise.
